% =====================================================================
%  skull-id-figures.tex
%  Figures for skull-id-guide.tex. Kept separate so they can be reused
%  (and so the schematics can be exported to SVG for a web version).
%
%  Two kinds of figure live here:
%
%   * SCHEMATICS, drawn in TikZ. Always available, no licence attached,
%     ours to relabel. These are deliberately diagrams: a beginner needs
%     the topology, and an honest map teaches that better than a mediocre
%     drawing of a real skull would.
%
%   * PLATES, clipped from the literature by docs/teaching/clip-plates.py
%     into docs/teaching/plates/ (git-ignored). Every one is wrapped in
%     \IfFileExists, so the guide compiles whether or not they are there.
%     Credit lines are set by the licences and are not optional.
%
%  Requires the colours defined in skull-id-guide.tex.
% =====================================================================

% ---------------------------------------------------------------------
%  Plate helper: include if clipped, otherwise leave a visible stub
% ---------------------------------------------------------------------
\newcommand{\plate}[4]{%
  % #1 = file stem in plates/  #2 = width  #3 = caption  #4 = credit
  \par\medskip
  \IfFileExists{plates/#1.png}{%
    \begin{center}\includegraphics[width=#2\linewidth]{plates/#1}\end{center}
    \vspace{-2pt}
    {\footnotesize #3\ {\color{stubgrey}\itshape #4}\par}
  }{%
    \begin{center}
      \fcolorbox{stubgrey}{white}{\parbox{0.86\linewidth}{\footnotesize
        \color{stubgrey}\textbf{Plate not clipped:} \texttt{plates/#1.png}.
        Run \texttt{python3 docs/teaching/clip-plates.py} with the PDFs in
        \texttt{papers/}.\\[3pt] #3\ \itshape #4}}
    \end{center}
  }%
  \par\medskip
}

% =====================================================================
%  FIGURE 1: skull-roof topology map
%  Tiles in their correct relative positions. What a beginner has to hold
%  is which bone touches which, so this is drawn as a map. Coordinates are
%  in centimetres so nothing collides when the text width changes.
% =====================================================================
\tikzset{
  bonetile/.style  = {draw=boneline, line width=0.5pt, rounded corners=1.5pt,
                      font=\scriptsize, align=center, inner sep=2.2pt,
                      minimum height=0.48cm},
  tstable/.style   = {bonetile, fill=stablecol},   % present, in position, ~always
  tlabile/.style   = {bonetile, fill=labilecol},   % lost or fused somewhere
  thinge/.style    = {bonetile, fill=hingecol},    % leaves the jaw in mammals
}

% \bonetile places a bone box at a named coordinate, rotated to sit roughly
% the way the bone sits. Contact lines are drawn BEFORE the boxes, so each one
% is occluded except where it crosses the gap between two neighbours -- which
% is exactly where an articulation is.
\tikzset{
  bonebox/.style   = {draw=boneline, line width=0.5pt, rounded corners=1.2pt,
                      align=center, inner xsep=1.6pt, inner ysep=1pt,
                      font=\tiny},
  bstable/.style   = {bonebox, fill=stablecol},
  blabile/.style   = {bonebox, fill=labilecol},
  bhinge/.style    = {bonebox, fill=hingecol},
  bdeep/.style     = {bonebox, fill=hingecol, densely dashed},
  bextra/.style    = {bonebox, fill=stubgrey!22},
  contact/.style   = {draw=stubgrey, line width=0.4pt},
  leader/.style    = {draw=stubgrey, line width=0.4pt},
  exlab/.style     = {font=\scriptsize, color=black, inner sep=1pt, align=left},
}

\newcommand{\figBoneMap}{%
  \par\medskip
  \noindent\begin{minipage}{\linewidth}
  \centering
  \begin{tikzpicture}[x=1cm, y=1cm]

    % ---------- bone centres -------------------------------------------
    \coordinate (pmx) at (0.80,2.28);   \coordinate (mx)  at (3.05,1.92);
    \coordinate (na)  at (2.45,3.62);   \coordinate (la)  at (4.30,2.78);
    \coordinate (prf) at (4.45,3.58);   \coordinate (fr)  at (5.40,4.10);
    \coordinate (pof) at (6.85,3.58);   \coordinate (po)  at (7.05,2.74);
    \coordinate (pa)  at (8.25,4.08);   \coordinate (ju)  at (6.40,2.02);
    \coordinate (sq)  at (9.35,3.18);   \coordinate (qj)  at (8.70,1.98);
    \coordinate (qu)  at (9.72,1.88);   \coordinate (st)  at (10.42,2.48);
    \coordinate (de)  at (3.35,0.88);   \coordinate (co)  at (6.30,1.44);
    \coordinate (sa)  at (7.50,1.00);   \coordinate (an)  at (7.20,0.44);
    \coordinate (sp)  at (4.90,0.46);   \coordinate (ar)  at (9.38,0.70);

    % ---------- articulations (drawn first, occluded by the boxes) ------
    \foreach \bxa/\bxb in {pmx/mx, pmx/na, na/fr, na/prf, na/la, mx/la, mx/ju,
                           la/prf, prf/fr, fr/pof, fr/pa, pof/po, pof/pa,
                           po/ju, po/sq, pa/sq, ju/qj, qj/qu, sq/qu, qu/ar,
                           de/sp, de/co, de/sa, de/an, sa/an, sa/ar, an/ar,
                           co/sa, sq/st}
      {\draw[contact] (\bxa) -- (\bxb);}

    % ---------- orbit ---------------------------------------------------
    \draw[draw=boneline, line width=0.6pt, densely dashed]
      (5.72,3.02) ellipse (0.80 and 0.62);
    \node[font=\scriptsize\itshape, color=boneline] at (5.72,3.02) {orbit};

    % ---------- upper skull --------------------------------------------
    \node[blabile, minimum width=0.92cm, minimum height=1.20cm, rotate=-18] at (pmx)
      {Pre-\\maxilla};
    \node[bstable, minimum width=3.6cm, minimum height=0.76cm, rotate=-4] at (mx) {Maxilla};
    \node[bstable, minimum width=2.8cm, minimum height=0.62cm, rotate=8]  at (na) {Nasal};
    \node[blabile, minimum width=0.82cm, minimum height=0.78cm] at (la) {};
    \node[blabile, minimum width=0.82cm, minimum height=0.70cm, rotate=20] at (prf) {};
    \node[bstable, minimum width=1.85cm, minimum height=0.64cm, rotate=3]  at (fr) {Frontal};
    \node[blabile, minimum width=0.84cm, minimum height=0.62cm, rotate=-18] at (pof) {};
    \node[blabile, minimum width=2.5cm, minimum height=0.60cm, rotate=9]  at (ju) {Jugal*};
    \node[blabile, minimum width=1.12cm, minimum height=1.02cm, rotate=-6] at (po)
      {Post-\\orbital*};
    \node[bstable, minimum width=1.95cm, minimum height=0.64cm, rotate=-5] at (pa) {Parietal};
    \node[bstable, minimum width=1.55cm, minimum height=0.92cm, rotate=-24] at (sq)
      {Squa-\\mosal};
    \node[blabile, minimum width=1.15cm, minimum height=0.56cm, rotate=-10] at (qj) {};
    \node[bhinge, minimum width=0.70cm, minimum height=1.55cm, rotate=7] at (qu) {};
    

    % ---------- lower jaw ------------------------------------------------
    \node[bstable, minimum width=4.5cm, minimum height=0.72cm, rotate=-3] at (de) {Dentary};
    \node[blabile, minimum width=0.98cm, minimum height=0.48cm, rotate=12] at (co)
      {Coronoid*};
    \node[blabile, minimum width=1.95cm, minimum height=0.56cm, rotate=3] at (sa)
      {Surangular*};
    \node[blabile, minimum width=2.05cm, minimum height=0.46cm, rotate=3] at (an)
      {Angular*};
    \node[blabile, minimum width=1.85cm, minimum height=0.42cm, rotate=-2] at (sp)
      {Splenial*};
    \node[bhinge, minimum width=0.78cm, minimum height=0.54cm, rotate=5] at (ar) {};
    \node[bdeep,  minimum width=0.95cm, minimum height=0.32cm, rotate=-14] at (st) {};

    % ---------- external labels: the small, crowded elements -------------
    % above the roof
    \draw[leader] (4.40,3.96) -- (4.15,4.95);
    \node[exlab, anchor=south] at (4.15,4.98) {Prefrontal*};
    \draw[leader] (6.88,3.92) -- (7.05,4.95);
    \node[exlab, anchor=south] at (7.05,4.98) {Postfrontal*};
    % in front
    \draw[leader] (3.95,2.92) -- (3.34,2.98);
    \node[exlab, anchor=east] at (3.30,2.98) {Lacrimal*};
    % behind
    \draw[leader] (10.86,2.60) -- (11.25,3.45);
    \node[exlab, anchor=west] at (11.22,3.62)
      {\textbf{Stapes} / columella\\[-1pt]{\tiny the hyomandibula, deep to the cheek}};
    \draw[leader] (10.10,1.95) -- (11.25,2.20);
    \node[exlab, anchor=west] at (11.25,2.20)
      {\textbf{Quadrate}\\[-1pt]{\tiny$\rightarrow$ incus}};
    \draw[leader] (9.78,0.70) -- (11.25,0.80);
    \node[exlab, anchor=west] at (11.25,0.80)
      {\textbf{Articular}\\[-1pt]{\tiny$\rightarrow$ malleus}};
    \draw[leader] (8.78,1.72) -- (8.58,-0.22);
    \node[exlab, anchor=north] at (8.58,-0.25) {Quadratojugal*};
    

    % ---------- the jaw joint -------------------------------------------
    

  \end{tikzpicture}

  \vspace{2pt}
  {\footnotesize\raggedright Generalised amniote, left lateral view. Boxes sit
   roughly where and how the bones sit, and the thin grey lines mark the usual
   contacts. The two pink elements meet at the jaw joint, and the stapes is
   dashed because it lies deep to the cheek.\par}
  \vspace{4pt}
  {\footnotesize\raggedright
   \tikz[baseline=-0.6ex]\node[bstable, minimum width=0.42cm,
     minimum height=0.28cm]{};\,\textbf{Stable}: present, and in this position,
   across essentially all bony jawed vertebrates.\quad
   \tikz[baseline=-0.6ex]\node[blabile, minimum width=0.42cm,
     minimum height=0.28cm]{};\,\textbf{Labile*}: lost or fused to a neighbour
   somewhere, and worth learning as absences.\quad
   \tikz[baseline=-0.6ex]\node[bhinge, minimum width=0.42cm,
     minimum height=0.28cm]{};\,\textbf{Hinge}: leaves the jaw entirely in
   mammals.\par}
  \end{minipage}
  \par\medskip
}

% =====================================================================
%  FIGURES 0a/0b: where the bones come from
%  Derivation trees for the neurocranium and the first two arches. Deliberately
%  uncoloured: figure 1 already spends colour on stability, and reusing the
%  lab's unit colours here would make pink and green mean two things each.
% =====================================================================
\tikzset{
  srcbox/.style  = {draw=boneline, line width=0.7pt, rounded corners=2pt,
                    fill=bonecol!70, font=\scriptsize\bfseries, align=center,
                    inner sep=3pt, minimum height=0.6cm},
  derbox/.style  = {draw=boneline, line width=0.4pt, rounded corners=2pt,
                    fill=bonecol!35, font=\tiny, align=center,
                    inner sep=2.5pt, minimum height=0.55cm},
  endbox/.style  = {derbox, fill=white, densely dashed},
  flow/.style    = {-{Stealth[length=4pt]}, draw=barcol, line width=0.6pt},
  qual/.style    = {font=\tiny\itshape, color=stubgrey, align=left},
}

\newcommand{\figNeurocranium}{%
  \par\smallskip
  \noindent\begin{minipage}{\linewidth}
  \centering
  \begin{tikzpicture}[x=1cm, y=1cm]
    \foreach \yy/\name in {4.35/Ethmoid, 2.95/Sphenoid, 1.55/Otic, 0.15/Occipital}
      {\node[srcbox, minimum width=1.8cm] at (0.9,\yy) {\name};
       \draw[flow] (1.85,\yy) -- (2.35,\yy);}

    % ethmoid
    \node[derbox] (e1) at (3.25,4.35) {Sphen-\\ethmoid};
    \node[derbox] (e2) at (4.75,4.35) {Mes-\\ethmoid};
    \node[derbox] (e3) at (6.15,4.35) {Turbinates};
    \node[derbox] (e4) at (7.75,4.35) {Cribriform\\plate};
    \node[derbox] (e5) at (9.35,4.35) {Ect-\\ethmoid};
    \node[qual, anchor=west] at (10.25,4.35)
      {Mostly stays cartilage. Sphenethmoid in frogs,\\
       cribriform plate in mammals, ectethmoid in \emph{Sphenodon}.};

    % sphenoid
    \node[derbox] (s1) at (3.35,2.95) {Basi-\\sphenoid};
    \node[derbox] (s2) at (4.95,2.95) {Pre-\\sphenoid};
    \node[derbox] (s3) at (6.65,2.95) {Latero-\\sphenoid};
    \node[derbox] (s4) at (8.35,2.95) {Orbito-\\sphenoid};
    \node[qual, anchor=west] at (9.35,2.95)
      {Basisphenoid under the midbrain everywhere. Presphenoid in some\\
       mammals, the other two in archosaurs. In most mammals all of them\\
       fuse into one \textbf{sphenoid}, its \emph{sella turcica} holding the pituitary.};

    % otic
    \node[derbox] (o1) at (3.25,1.55) {Prootic};
    \node[derbox] (o2) at (4.75,1.55) {Opisthotic};
    \node[derbox] (o3) at (6.15,1.55) {Epiotic};
    \node[qual, anchor=west] at (7.05,1.55)
      {Opisthotic fuses to the exoccipital in frogs and most amniotes; all three\\
       fuse in birds (the \textbf{periotic}) and in mammals (the \textbf{petrosal}).};

    % occipital
    \node[derbox] (c1) at (3.35,0.15) {Basi-\\occipital};
    \node[derbox] (c2) at (5.05,0.15) {Exoccipitals\\(paired)};
    \node[derbox] (c3) at (6.85,0.15) {Supra-\\occipital};
    \node[qual, anchor=west] at (7.85,0.15)
      {Ring the \emph{foramen magnum}. The exoccipitals carry the condyles,\\
       and salamanders keep only those two.};
  \end{tikzpicture}

  \vspace{2pt}
  {\footnotesize\raggedright The four neurocranial ossification centers (anterior to posterior) \& their derivatives.\par}
  \end{minipage}
  \par\smallskip
}

\newcommand{\figArches}{%
  \par\smallskip
  \noindent\begin{minipage}{\linewidth}
  \centering
  \begin{tikzpicture}[x=1cm, y=1cm]
    % ---------------- arch 1
    \node[font=\scriptsize\bfseries, color=headcol, anchor=west] at (0,3.75)
      {Arch 1 (mandibular)};
    \node[srcbox, minimum width=2.3cm] (pq) at (1.35,2.95) {Palato-\\quadrate};
    \node[srcbox, minimum width=2.3cm] (mc) at (1.35,1.45) {Meckel's\\cartilage};

    \node[derbox, minimum width=1.5cm] (qd) at (4.45,3.35) {Quadrate};
    \node[derbox, minimum width=1.5cm] (ep) at (4.45,2.45) {Epi-\\pterygoid};
    \node[derbox, minimum width=1.5cm] (ar) at (4.45,1.45) {Articular};

    \node[endbox, minimum width=1.9cm] (ic) at (7.30,3.35) {Incus\\{\tiny mammals}};
    \node[endbox, minimum width=1.9cm] (al) at (7.30,2.45) {Alisphenoid\\{\tiny mammals, probable}};
    \node[endbox, minimum width=1.9cm] (ma) at (7.30,1.45) {Malleus\\{\tiny mammals}};

    \foreach \a/\b in {pq/qd, pq/ep, mc/ar} {\draw[flow] (\a) -- (\b);}
    \foreach \a/\b in {qd/ic, ep/al, ar/ma} {\draw[flow] (\a) -- (\b);}

    \node[qual, anchor=west] at (8.60,2.40)
      {The dentary, splenial, angular, surangular and\\
       coronoid that sheath Meckel's cartilage are\\
       \textbf{dermal} bone, not arch 1.};

    % ---------------- arch 2
    \draw[draw=stubgrey, line width=0.3pt] (0,0.85) -- (16.2,0.85);
    \node[font=\scriptsize\bfseries, color=headcol, anchor=west] at (0,0.42)
      {Arch 2 (hyoid)};
    \node[srcbox, minimum width=2.3cm] (hm) at (1.35,-0.55) {Hyoman-\\dibula};
    \node[srcbox, minimum width=2.3cm] (ch) at (1.35,-1.85) {Cerato-\\hyal};

    \node[derbox, minimum width=2.1cm] (cl) at (4.70,-0.55) {Columella};
    \node[endbox, minimum width=1.9cm] (st) at (7.40,-0.55) {Stapes\\{\tiny mammals}};
    \node[derbox, minimum width=4.8cm] (sty) at (6.05,-1.85)
      {Styloid process, lesser horn of the hyoid,\\stylohyoid ligament};

    \foreach \a/\b in {hm/cl, cl/st, ch/sty} {\draw[flow] (\a) -- (\b);}
    \node[qual, anchor=west] at (8.60,-1.85)
      {Arches 3--7 give the rest of the hyoid\\and the laryngeal cartilages.};
  \end{tikzpicture}

  \vspace{2pt}
  {\footnotesize\raggedright The first two pharyngeal arches.\par}
  \end{minipage}
  \par\smallskip
}

% =====================================================================
%  FIGURE 1b: the same stylised skull across four groups
%  Same coordinates as figure 1, carrying the bone abbreviations. A bone that
%  a group lacks is simply not drawn, so the silhouette of what is missing is
%  the whole point. Codes are baked into the \bXX macros so the panels and the
%  caption key cannot drift apart.
% =====================================================================
% #1 x  #2 y  #3 w  #4 h  #5 rotation  #6 style  #7 abbreviation
\newcommand{\bn}[7]{\node[#6, minimum width=#3cm, minimum height=#4cm,
  rotate=#5, inner xsep=1pt, inner ysep=0.5pt] at (#1,#2) {#7};}

% every element of the generalised skull, as its own macro
\newcommand{\bPMX}[1]{\bn{0.80}{2.28}{0.92}{1.20}{-18}{#1}{PM}}
\newcommand{\bMX} [1]{\bn{3.05}{1.92}{3.60}{0.76}{-4}{#1}{MX}}
\newcommand{\bNA} [1]{\bn{2.45}{3.62}{2.80}{0.62}{8}{#1}{NA}}
\newcommand{\bLA} [1]{\bn{4.30}{2.78}{0.82}{0.78}{0}{#1}{LA}}
\newcommand{\bPRF}[1]{\bn{4.45}{3.58}{0.82}{0.70}{20}{#1}{PRF}}
\newcommand{\bFR} [1]{\bn{5.40}{4.10}{1.85}{0.64}{3}{#1}{FR}}
\newcommand{\bPOF}[1]{\bn{6.85}{3.58}{0.84}{0.62}{-18}{#1}{POF}}
\newcommand{\bPO} [1]{\bn{7.05}{2.74}{1.12}{1.02}{-6}{#1}{PO}}
\newcommand{\bPA} [1]{\bn{8.25}{4.08}{1.95}{0.64}{-5}{#1}{PA}}
\newcommand{\bJU} [1]{\bn{6.40}{2.02}{2.50}{0.60}{9}{#1}{JU}}
\newcommand{\bSQ} [1]{\bn{9.35}{3.18}{1.55}{0.92}{-24}{#1}{SQ}}
\newcommand{\bQJ} [1]{\bn{8.70}{1.98}{1.15}{0.56}{-10}{#1}{QJ}}
\newcommand{\bQU} [1]{\bn{9.72}{1.88}{0.70}{1.55}{7}{#1}{QU}}
\newcommand{\bST} [1]{\bn{10.98}{2.62}{0.95}{0.32}{-14}{#1}{ST}}
\newcommand{\bDE} [1]{\bn{3.35}{0.88}{4.50}{0.72}{-3}{#1}{DE}}
\newcommand{\bCO} [1]{\bn{6.30}{1.44}{0.98}{0.48}{12}{#1}{CO}}
\newcommand{\bSA} [1]{\bn{7.50}{1.00}{1.95}{0.56}{3}{#1}{SA}}
\newcommand{\bAN} [1]{\bn{7.20}{0.44}{2.05}{0.46}{3}{#1}{AN}}
\newcommand{\bSP} [1]{\bn{4.90}{0.46}{1.85}{0.42}{-2}{#1}{SP}}
\newcommand{\bAR} [1]{\bn{9.38}{0.70}{0.78}{0.54}{5}{#1}{AR}}
\newcommand{\bORB}{\draw[draw=boneline, line width=0.5pt, densely dashed]
  (5.72,3.02) ellipse (0.80 and 0.62);}

% \minipanel{caption}{body}
\newcommand{\minipanel}[2]{%
  \begin{minipage}[t]{0.48\linewidth}
    \centering
    \begin{tikzpicture}[scale=0.55, transform shape, x=1cm, y=1cm,
        every node/.append style={font=\large\bfseries}]
      #2
    \end{tikzpicture}\\[1pt]
    {\raggedright\scriptsize #1\par}
  \end{minipage}%
}

\newcommand{\figSkullVariants}{%
  \par\medskip
  \noindent
  \minipanel{\textbf{Actinopterygian.} Generalized skull retains everything with some bonus bones. Many, \textit{many} actinopterygians reduce and alter their skull dramatically. The base, ``generic'' skull has well over a hundred elements.}{%
      \bORB \bPMX{bstable}\bMX{bstable}\bNA{bstable}\bLA{blabile}\bPRF{blabile}
      \bFR{bstable}\bPOF{blabile}\bJU{blabile}\bPO{blabile}\bPA{bstable}
      \bSQ{bstable}\bQJ{blabile}\bQU{bhinge}\bDE{bstable}\bCO{blabile}
      \bSA{blabile}\bAN{blabile}\bSP{blabile}\bAR{bhinge}
      \bn{10.95}{2.20}{1.20}{2.00}{-5}{bextra}{POP}
      \bn{12.35}{3.00}{1.50}{2.20}{-6}{bextra}{OP}
      \bn{12.45}{1.35}{1.30}{1.20}{-8}{bextra}{SOP}
      \bn{11.30}{0.35}{2.20}{0.50}{-12}{bextra}{BR}}
  \hfill
  \minipanel{\textbf{Salamander.} Jugal, quadratojugal, postorbital, postfrontal
    and lacrimal all lost.}{%
      \bORB \bPMX{bstable}\bMX{bstable}\bNA{bstable}\bPRF{blabile}
      \bFR{bstable}\bPA{bstable}\bSQ{bstable}\bQU{bhinge}\bST{bdeep}
      \bDE{bstable}\bCO{blabile}\bSA{blabile}\bAN{blabile}\bSP{blabile}
      \bAR{bhinge}}
  \par\medskip
  \noindent
  \minipanel{\textbf{Lizard.} The quadratojugal is lost, resulting in no lower temporal bar and a situation where the quadrate can swing freely on the squamosal.}{%
      \bORB \bPMX{bstable}\bMX{bstable}\bNA{bstable}\bLA{blabile}\bPRF{blabile}
      \bFR{bstable}\bPOF{blabile}\bJU{blabile}\bPO{blabile}\bPA{bstable}
      \bSQ{bstable}\bQU{bhinge}\bST{bdeep}\bDE{bstable}\bCO{blabile}
      \bSA{blabile}\bAN{blabile}\bSP{blabile}\bAR{bhinge}}
  \hfill
  \minipanel{\textbf{Mammal.} Lower jaw is only the mandible, with the post-dentary bones moving or fusing. The jaw joint is now dentary to squamosal directly. Postorbital, postfrontal, prefrontal and quadratojugal all lost.}{%
      \bORB \bPMX{bstable}\bMX{bstable}\bNA{bstable}\bLA{blabile}
      \bFR{bstable}\bJU{blabile}\bPA{bstable}\bSQ{bstable}
      \bDE{bstable}
      \bn{10.80}{3.20}{0.80}{0.55}{8}{bhinge}{IN}
      \bn{12.00}{3.20}{0.85}{0.55}{-4}{bhinge}{MA}
      \bn{13.20}{3.20}{0.80}{0.45}{-10}{bdeep}{ST}}
  \par\medskip
  {\footnotesize\raggedright Generalized skull diagrams for major clades. Colors as in the map above; grey boxes show bones actinopterygians have that aren't part of the tetrapod skull.
   AN, angular; AR, articular; BR, branchiostegals; CO, coronoid; DE, dentary;
   FR, frontal; IN, incus; JU, jugal; LA, lacrimal; MA, malleus; MX, maxilla;
   NA, nasal; OP, opercular; PA, parietal; PM, premaxilla; PO, postorbital;
   POF, postfrontal; POP, preopercular; PRF, prefrontal; QJ, quadratojugal; QU, quadrate; SA,
   surangular; SOP, subopercular; SP, splenial; SQ, squamosal; ST, stapes.\par}
  \par\medskip
}

% =====================================================================
%  FIGURE 2: the temporal series
% =====================================================================
\tikzset{
  skullfill/.style = {fill=bonecol, draw=boneline, line width=0.6pt},
  openfill/.style  = {fill=white, draw=boneline, line width=0.5pt},
  fenfill/.style   = {fill=fencol, draw=boneline, line width=0.5pt},
  barbold/.style   = {draw=barcol, line width=1.7pt, line cap=round},
  cuthint/.style   = {draw=barcol, line width=0.8pt,
                      dash pattern=on 2pt off 1.7pt},
}

% Generic amniote lateral silhouette. Snout left, occiput right.
\newcommand{\skulloutline}{%
  \path[skullfill]
    (0,0.60)
      .. controls (0.55,1.00) and (1.35,1.22) .. (2.15,1.26)
      .. controls (2.85,1.29) and (3.32,1.16) .. (3.42,0.90)
    -- (3.38,0.34) -- (2.60,0.16) -- (0.06,0.14) -- cycle;
}
% Orbit and naris never move. That is the point: only the back changes.
\newcommand{\skullorbit}{%
  \path[openfill] (1.85,0.76) circle (0.30);
  \path[openfill] (0.30,0.66) ellipse (0.12 and 0.09);
}

% \fenpanel{title}{note}{extra tikz}
\newcommand{\fenpanel}[3]{%
  \begin{minipage}[t]{0.235\linewidth}
    \centering
    \begin{tikzpicture}[scale=1.02]
      \skulloutline #3 \skullorbit
    \end{tikzpicture}\\[3pt]
    {\footnotesize\bfseries #1}\\[2pt]
    \begin{minipage}{\linewidth}
      \raggedright\scriptsize\color{stubgrey}\setlength{\parskip}{0pt}#2
    \end{minipage}
  \end{minipage}%
}

\newcommand{\figFenestration}{%
  \par\medskip
  \noindent
  \begin{tabular}{@{}c@{\hspace{0.012\linewidth}}c@{\hspace{0.012\linewidth}}%
                   c@{\hspace{0.012\linewidth}}c@{}}
  \fenpanel{Anapsid}{No opening at all. Stem amniotes, and the starting
    condition for everything else in this figure.}{}
  &
  \fenpanel{Synapsid}{One opening, and it is \emph{low}. The postorbital
    and squamosal meet above it to close the top.}{%
      \path[fenfill] (2.72,0.55) ellipse (0.33 and 0.20);
      \draw[barbold] (2.32,0.92) -- (3.20,0.88);}
  &
  \fenpanel{Diapsid}{Two openings, one above the other, separated by the
    postorbital--squamosal bar. Tuatara, crocodilians.}{%
      \path[fenfill] (2.78,1.02) ellipse (0.32 and 0.14);
      \path[fenfill] (2.74,0.55) ellipse (0.31 and 0.19);
      \draw[barbold] (2.34,0.82) -- (3.22,0.78);
      \draw[barbold] (2.36,0.28) -- (3.30,0.30);}
  &
  \fenpanel{Lizard}{Upper opening and its bar kept. The \emph{lower bar}
    is gone, so that opening runs out to the margin.}{%
      \path[fenfill] (2.78,1.02) ellipse (0.32 and 0.14);
      \path[fenfill] (2.28,0.62) -- (3.36,0.58) -- (3.38,0.22)
                     -- (2.52,0.16) -- cycle;
      \draw[barbold] (2.34,0.82) -- (3.22,0.78);
      \draw[cuthint] (2.42,0.24) -- (3.34,0.26);}
  \\[10pt]
  \fenpanel{Turtle}{No fenestra anywhere, but the cheek is deeply
    \emph{emarginated}, notched inwards from the free rear edge.}{%
      \path[fill=white]
        (3.46,1.02) .. controls (2.88,0.86) and (2.74,0.46) ..
        (3.04,0.12) -- (3.54,0.12) -- (3.54,1.02) -- cycle;
      \draw[draw=boneline, line width=0.6pt]
        (3.42,0.92) .. controls (2.88,0.84) and (2.76,0.46) ..
        (3.05,0.14);}
  &
  \fenpanel{Bird}{Postorbital bar lost, so orbit and temporal region run
    together into one great opening. The jugal bar is a bare rod.}{%
      \path[fill=white, draw=boneline, line width=0.5pt]
        (2.12,1.10) .. controls (2.70,1.12) and (3.06,0.96) ..
        (3.14,0.64) .. controls (3.00,0.48) and (2.50,0.50) ..
        (2.14,0.56) -- cycle;
      \draw[barbold] (1.30,0.28) -- (3.20,0.33);}
  &
  \fenpanel{Mammal}{One opening, hugely enlarged; postorbital bar usually
    gone, and the \emph{zygomatic arch} forms the floor of the fossa.}{%
      \path[fill=white, draw=boneline, line width=0.5pt]
        (2.08,1.16) .. controls (2.80,1.18) and (3.24,1.00) ..
        (3.30,0.62) .. controls (3.00,0.50) and (2.40,0.52) ..
        (2.10,0.58) -- cycle;
      \draw[barbold] (1.55,0.42) .. controls (2.40,0.28) and (3.00,0.30)
        .. (3.32,0.42);}
  &
  \begin{minipage}[t]{0.235\linewidth}
    \raggedright\footnotesize
    \textbf{Reading the panels}\\[3pt]
    {\scriptsize
     \tikz[baseline=-0.5ex]\draw[barbold](0,0)--(0.5,0); \; a \textbf{bar}:
     bone you can trace with a finger.\\[3pt]
     \tikz[baseline=-0.5ex]\path[fenfill](0,0) ellipse (0.26 and 0.13);
     \; a \textbf{fenestra}: a hole ringed by bone.\\[3pt]
     \tikz[baseline=-0.5ex]\draw[cuthint](0,0)--(0.5,0); \; where a bar has
     been \textbf{lost}.\\[4pt]
     Orbit and nostril are drawn in the same place every time. Only the
     temporal region changes.}
  \end{minipage}
  \end{tabular}
  \par\medskip
}

% =====================================================================
%  FIGURE 3: the jaw joint becomes the middle ear
% =====================================================================
\newcommand{\figOssicles}{%
  \par\medskip
  \begin{center}
  \begin{tikzpicture}[x=1cm, y=1cm,
      oldbox/.style={draw=boneline, line width=0.5pt, rounded corners=2.5pt,
                     fill=bonecol, font=\small, align=center,
                     minimum width=2.9cm, minimum height=0.68cm},
      newbox/.style={draw=boneline, line width=0.5pt, rounded corners=2.5pt,
                     fill=hingecol, font=\small, align=center,
                     minimum width=2.9cm, minimum height=0.68cm},
      arr/.style={-{Stealth[length=5pt]}, line width=0.9pt, color=barcol},
      viaz/.style={font=\scriptsize\itshape, color=stubgrey, align=center,
                   fill=white, inner sep=1.5pt}]

    \node[font=\footnotesize\bfseries, color=headcol] at (0,3.45)
      {Non-mammal: in the jaw};
    \node[font=\footnotesize\bfseries, color=headcol] at (8.2,3.45)
      {Mammal: in the middle ear};

    \foreach \y/\old/\new/\via in {
        2.7/Hyomandibula/Stapes/{already a sound conductor},
        1.8/Quadrate/Incus/{upper half of the jaw joint},
        0.9/Articular/Malleus/{lower half of the jaw joint},
        0.0/Angular/{Tympanic bulla}/{fuses to the skull}}
    {
      \node[oldbox] (o) at (0,\y) {\old};
      \node[newbox] (n) at (8.2,\y) {\new};
      \draw[arr] (1.48,\y) -- (6.72,\y);
      \node[viaz] at (4.1,\y) {\via};
    }

    \node[draw=barcol, line width=0.7pt, rounded corners=2.5pt,
          fill=fencol, font=\footnotesize, align=center, text width=10.4cm,
          inner sep=5pt]
      at (4.1,-1.25)
      {With the hinge elements gone, the \textbf{dentary} articulates directly with the \textbf{squamosal}.};
  \end{tikzpicture}
  \end{center}
  \par\medskip
}
